\documentclass[11pt]{article}
\usepackage[margin=1in]{geometry}
\usepackage{amsmath,amssymb,amsthm,mathtools}
\usepackage{booktabs}
\usepackage{graphicx}
\usepackage{enumitem}
\usepackage[hidelinks]{hyperref}
\hypersetup{pdftitle={A marching-compatible separable test norm}}

\newtheorem{proposition}{Proposition}
\theoremstyle{remark}
\newtheorem{remark}[proposition]{Remark}

\newcommand{\bbeta}{\boldsymbol{\beta}}
\newcommand{\bn}{\boldsymbol{n}}
\newcommand{\bv}{\boldsymbol{v}}
\newcommand{\T}{\mathsf{T}}
\newcommand{\R}{\mathbb{R}}
\DeclareMathOperator{\diag}{diag}

\title{A marching-compatible separable test norm\\[2pt]\large proposal, for the tensor-train and conventional planar MITL solvers}
\author{}
\date{November 20, 2026}

\begin{document}
\maketitle

\begin{abstract}
The transport-sweep experiments point at the test norm, rather than the solver, as the source of two of the three
degradations we found.  The manuscript ``The DPG Method for the Convection--Reaction Problem, Revisited''
(Demkowicz, Roberts, Mu\~noz-Matute) recalls an inner product that replaces the $L^2$ term by an \emph{outflow-boundary}
term, and with which the ideal DPG method coincides with a marching (sweep) scheme in 1D and on a class of 2D meshes.
These notes propose a direction-separable version of that norm.  Its element Gram matrix is still a Kronecker sum, so
its inverse is still exact by fast diagonalization.  What changes is that the one-dimensional eigenbases become
cell-dependent: they depend on one or two scalars per direction and cell.  The proposal is untested.  It changes the
discretization, so it needs the accuracy and coupled-stability tests as well as solver tests.
\end{abstract}

%======================================================================
\section{Background: the current separable norm and what the sweep found}
%======================================================================

\paragraph{Current norm (TT ``graph$\oplus$'').}  Coordinates $\mathbf z=(x,t,u_x,u_z)$, $n=4$.  On a box cell $K$ with
widths $h_{K,d}$:
\[
  \|w\|_{V(K)}^2 = \alpha_K\|w\|_K^2 + \sum_{d=1}^n s_{K,d}^2\,\|\partial_d w\|_K^2\;\bigl[-\,c_K\bigl(\textstyle\int_K w\bigr)^2\bigr],
\]
where:
\begin{itemize}[nosep]
  \item $s_K=(v_{x,\max}(K),\,1,\,a_x,\,a_z)$ are the derivative scales, with $a=(1.8,0.65)$ fixed;
  \item $\alpha_K=\alpha+C\bigl[(v_{x,\max}/h_x)^2+1/h_t^2+\sum_u(a_u/h_u)^2\bigr]$ is the transport-scaled $L^2$
        weight, with $\alpha=10^2$ and $C=1.6$;
  \item the bracketed constant-mode term is optional.
\end{itemize}
On the reference cell $[-1,1]^n$, with Jacobian $J_K=\prod_d h_{K,d}/2$, the Gram matrix is
\[
  G_K = J_K\Bigl(\alpha_K\,\widehat M^{\otimes n}+\sum_d \xi_{K,d}\,\widehat K_{(d)}\Bigr),\qquad \xi_{K,d}=s_{K,d}^2\,(2/h_{K,d})^2 .
\]
Here $\widehat M$ and $\widehat K$ are the 1D mass and stiffness matrices of the test basis, and $\widehat K_{(d)}$ has
$\widehat K$ in slot $d$ and $\widehat M$ elsewhere.  One 1D generalized eigenbasis $V$, with $\widehat KV=\widehat MV\Lambda$,
diagonalizes every cell: $G_K^{-1}=V^{\otimes n}\diag\bigl(1/(J_K(\alpha_K+\sum_d\xi_{K,d}\lambda_{j_d}))\bigr)(V^\T)^{\otimes n}$.

\paragraph{The sweep and what limits it.}  The notes \texttt{static-condensation-trace-structure.pdf} split each
cell's faces into inflow and outflow.  For that split they show that the DPG operator is exactly
\[
  S = R^\T W R + Q,\qquad Q\succeq 0,
\]
where $R$ is the sweep (each cell's local DPG solve from its inflow data, in upwind order) and $Q$ collects the cells'
inflow Schur complements.  The sweep preconditioner $P=(R^\T WR)^{-1}$ then has $\kappa(PS)\le 1+\rho$, with
$\rho=\sup\, x^\T Qx/x^\T R^\T WRx$.  $Q$ vanishes in the continuous limit, so $\rho$ measures how far the
\emph{discrete} local problems are from causal ``inflow determines outflow'' behaviour.  Measured ($p=2$):
\begin{center}
\resizebox{\textwidth}{!}{%
\begin{tabular}{@{}llll@{}}
\toprule
case & what dominated $\rho$ & test & effect\\
\midrule
fields zero, 16 graded $x$-cells (thin cathode cells) & t-face fluxes in the thinnest cells & $\alpha_K$ fixed at $10^2$
  & sweep $\kappa$: 6036 $\to$ 6.2\\
fields on, production L0 (880 cells) & $u_x$-face fluxes, cathode cells, $|\beta_{u_x}|\approx4$--6 $>$ $a_x=1.8$
  & $a=(6.11,4.15)$ & sweep $\kappa$: 109 $\to$ 43\\
 & & & two-level: 15 $\to$ 8 iterations\\
fields on, production L0 & cycles from $\bv\times\mathbf B$ (one strongly connected component) & weighted lagging
  & sweep $\kappa$: 769 $\to$ 109\\
\bottomrule
\end{tabular}}
\end{center}
The first two are properties of the norm:
\begin{enumerate}[nosep]
  \item \textbf{$\alpha_K$ dominating the time derivative.}  In a thin cathode cell,
        $\alpha_K\approx C(v_x/h_x)^2$, while the time-derivative weight is $1/h_t^2$.  Relative to the $L^2$ term,
        $\partial_t$ is down by about $1/(C\nu^2)$, with $\nu=|v_x|h_t/h_x$ the cell Courant number.  An
        $L^2$-dominated local problem is close to a least-squares fit that is symmetric in $t$.  It does not ``know''
        that the bottom face is inflow and the top face is outflow.  (Production meshes have $\nu\approx13$ at the
        cathode, where $\alpha_K\approx10^5$ against $C/h_t^2=640$ at L1.)
  \item \textbf{Domination failure} where $s_{K,d}<|\beta_d|$ (the fixed acceleration scales).
\end{enumerate}

%======================================================================
\section{The manuscript's inner product}
%======================================================================

For $\mathbf b\cdot\nabla u+cu=f$, the manuscript recalls, from its reference [12], the element inner product
\begin{equation}\label{eq:11}
  (v,\delta v)_{V(K)} = \int_K \operatorname{div}(\mathbf b v)\operatorname{div}(\mathbf b\,\delta v) + \int_{\partial K_+} v\,\delta v .
\end{equation}
\begin{itemize}[nosep]
  \item With \eqref{eq:11}, the optimal test functions are explicit.  The ideal DPG method then delivers exactly the
        solution of the marching scheme, which solves element by element from inflow to outflow, in 1D and for a class of
        2D meshes.  For general 2D meshes the two differ slightly.
  \item \eqref{eq:11} is \emph{not} uniformly equivalent in $h$ to the adjoint graph norm.  This is footnote 2, citing
        Broersen, Dahmen and Stevenson, who use a corrected version.
  \item Their discrete stability argument (Section 6) is itself a layer-by-layer sweep.  The global constant is a
        product of element constants $\beta_h^{-2}$ along characteristic chains, and it needs $\beta_h^{-2}\le1+Ch$.
        $\beta_h$ degenerates for elements with more than one outflow face when the test space is a plain enriched
        polynomial space.  Their fix is a composite test space, split along a streamline so that each piece has one
        outflow face.
  \item Trace stability degenerates on faces nearly parallel to the streamlines, and it is much weaker at $p=4$ with
        $\Delta p=1$ (Table 5).  $\Delta p=2$ restores $h$-independence.
\end{itemize}
In our notation, marching equivalence means $Q=0$: the sweep is then an exact solver.  The hope is that a separable
version of \eqref{eq:11} keeps $Q$, and therefore $\rho$, small.

%======================================================================
\section{The proposal: a separable, marching-compatible norm}
%======================================================================

\subsection{Definition}

For each cell $K$ and direction $d$, let $F^+_{K,d}$ be the outflow face of $K$ normal to $\mathbf e_d$: the upper face
if $\beta_d>0$ on $K$ and the lower face if $\beta_d<0$.  Let $\mu_{K,d}\ge0$ be a weight for it.  Define
\begin{equation}\label{eq:norm}
  \|w\|_{V(K)}^2 = \sum_{d=1}^n s_{K,d}^2\,\|\partial_d w\|_K^2 + \sum_{d=1}^n \mu_{K,d}\int_{F^+_{K,d}} w^2
  \;+\;\alpha_K'\|w\|_K^2 \;\bigl[-\,c_K\bigl(\textstyle\int_K w\bigr)^2\bigr].
\end{equation}
\begin{itemize}[nosep]
  \item The first two sums are the separable surrogate of \eqref{eq:11}.  The derivative term dominates
        $\|\bbeta\cdot\nabla w\|^2/n$ when $s_{K,d}\ge|\beta_d|$ on $K$, and $\operatorname{div}\bbeta=0$.  The
        outflow term is \eqref{eq:11}'s boundary term.  The natural weight is $\mu_{K,d}\approx|\beta_d|$ on the face,
        since $\int_{\partial K_+}(\bbeta\cdot\bn)\,w^2$ is the weighted form used in the manuscript's trace lemmas.
  \item $\alpha_K'\ge0$ is an optional, \emph{small} $L^2$ term.  With $\alpha_K'=0$, \eqref{eq:norm} is still a norm
        on the test space: a zero derivative term forces $w$ to be constant, and the outflow term then forces the
        constant to be zero.
  \item The constant-mode correction stays a rank-one update, handled by Sherman--Morrison as now.
\end{itemize}

\subsection{The Gram matrix is still a Kronecker sum}

On the reference cell, the face normal to $d$ at endpoint $\pm1$ has $\int_F w^2 = J_K\,(2/h_{K,d})\int_{\widehat F}\hat w^2$.
Its 1D factor in direction $d$ is the endpoint mass matrix
$E_\pm=\boldsymbol\phi(\pm1)\boldsymbol\phi(\pm1)^\T$ (rank one), with $\widehat M$ in the other slots.  Splitting
$\alpha_K'$ evenly over the directions,
\begin{equation}\label{eq:gram}
  G_K = J_K\sum_{d=1}^n \widehat M\otimes\cdots\otimes\underbrace{A_{K,d}}_{\text{slot }d}\otimes\cdots\otimes\widehat M,
  \qquad A_{K,d} = \xi_{K,d}\,\widehat K + \eta_{K,d}\,E_{\pm} + \tfrac{\alpha_K'}{n}\,\widehat M,
\end{equation}
with $\xi_{K,d}=s_{K,d}^2(2/h_{K,d})^2$, $\eta_{K,d}=\mu_{K,d}\,(2/h_{K,d})$, and $\pm$ the outflow side.

\begin{proposition}[exact inverse; proved]
Let $A_{K,d}V_{K,d}=\widehat M V_{K,d}\Lambda_{K,d}$ with $V_{K,d}^\T\widehat MV_{K,d}=I$, one $m\times m$ generalized
eigenproblem per direction, where $m=p+\Delta p+1$.  Then
\[
  G_K^{-1} = \Bigl(\bigotimes_d V_{K,d}\Bigr)\,\diag\Bigl(\frac{1}{J_K\sum_d\lambda_{K,d,j_d}}\Bigr)\Bigl(\bigotimes_d V_{K,d}\Bigr)^{\!\T}.
\]
\end{proposition}
\begin{proof}
Each term of \eqref{eq:gram} is diagonalized by $\bigotimes_dV_{K,d}$, because $V_{K,d'}^\T\widehat MV_{K,d'}=I$ in
every slot $d'\ne d$ and $V_{K,d}^\T A_{K,d}V_{K,d}=\Lambda_{K,d}$ in slot $d$.
\end{proof}

\paragraph{What changes in practice.}
\begin{itemize}[nosep]
  \item The 1D eigenbasis is no longer shared by all cells.  $V_{K,d}$ depends on the ratio
        $\eta_{K,d}/\xi_{K,d}$ (and $\alpha'_K/\xi_{K,d}$), on the outflow side, and on nothing else.  Per cell this means
        $n$ eigenproblems of size $m\le9$ (negligible), or a table over a few ratios.
  \item The sum-factorized application is unchanged: mode products, now with per-cell $V_{K,d}$.
  \item The ``shared-eigenbasis coordinates'' used by the matrix-free diagonal and block computations become
        per-cell coordinates.  They are still cheap.
\end{itemize}

\paragraph{Field dependence (relevant for TT reuse).}  The outflow sides in $t$ and $x$ do not depend on the fields:
$\beta_t=1$ (always the upper face), and $\beta_x=v_x(\mathbf u)$ has the sign of $u_x$, with $u_x=0$ a mesh vertex.
Only the $u$-directions' outflow sides and weights depend on the fields, through $a_u=(q/m)(E+v\times B)_u$.  So:
\begin{itemize}[nosep]
  \item \textbf{Variant A, field-independent:} outflow terms in $t$ and $x$ only.  Keep the $u$-directions as now:
        derivative scales $s_{K,u}$ (preferably the planned per-velocity-cell $a_u(K)$), plus a small $L^2$ term if
        wanted.  The element inverse stays reusable across substeps.
  \item \textbf{Variant B, full:} outflow terms in every direction, with $F^+_{K,u}$ from the sign of the current
        $a_u$ on $K$.  Where $a_u$ changes sign within a cell, use both faces with weights $|a_u|$ restricted to each
        sign, or fall back to variant A in that direction.  The inverses then have to be rebuilt when the signs
        change, i.e.\ at slab boundaries or on change.
\end{itemize}
Variant A already addresses degradation 1 (the $t$--$x$ causality in the thin cathode cells).  Per-cell $a_u(K)$
addresses degradation 2.

%======================================================================
\section{Why it should help, and the caveats}
%======================================================================

\paragraph{Expected effect on the solver (hypothesis).}  The degradations we measured came from local problems that
are not causal: an $L^2$-dominated norm, or one that under-weights a transport direction.  \eqref{eq:norm} penalizes
only derivatives and outflow traces, which is the structure under which DPG reduces to marching.  So the local DPG
solves should behave like ``inflow determines outflow'', $\Sigma_K$ and $Q$ should shrink, and $P$ should approach
$S^{-1}$.  The same structure should improve any preconditioner that relies on local, causal solves.  This is
plausible, not proved.  Even for \eqref{eq:11}, marching equivalence is known only in 1D and on special 2D meshes.
Our boxes have several outflow faces per cell, which is the manuscript's degenerate case.

\paragraph{Caveats.}
\begin{enumerate}[nosep]
  \item \textbf{It changes the discretization.}  The norm sets the optimal test functions and therefore the solution.
        \eqref{eq:11} is not uniformly equivalent to the graph norm in $h$, which affects robustness of the error
        estimates.  The separable surrogate adds its usual domination factor (at most $n$ when $s_{K,d}\ge|\beta_d|$).
  \item \textbf{Coupled stability.}  The current $\alpha_K$ was tuned to remove growing modes:
        \begin{itemize}[nosep]
          \item a fixed $\alpha=4\cdot10^2$ grew a negative-density mode in the cathode cells;
          \item $\max(10^2, 1.6/h_x^2)$ grew a mode in the anode cells;
          \item the $(a_u/h_u)^2$ terms removed both.
        \end{itemize}
        A norm with little or no $L^2$ term has to pass those tests again.
  \item \textbf{Several outflow faces.}  The manuscript's stability argument needs composite (sub-grid) test spaces
        for elements with more than one outflow face.  A tensor sub-grid test space, piecewise polynomial in each 1D
        factor, keeps \eqref{eq:gram}'s Kronecker structure.  It only approximates the streamline split, because
        $v_x$ varies with $\mathbf u$ inside a cell.
  \item \textbf{Characteristic faces} ($\beta_d\approx0$ on a face, e.g.\ $x$-faces near $u_x=0$).  $\mu_{K,d}$ should
        vanish with $|\beta_d|$ there, and the derivative scale $s_{K,d}$ keeps the norm definite.
\end{enumerate}

%======================================================================
\section{Suggested tests}
%======================================================================

\begin{enumerate}[nosep]
  \item \textbf{Implementation}: a \texttt{SeparableTensorIP} option for \eqref{eq:norm}, with per-cell 1D
        eigenbases.  The current shared-basis path stays the default.
  \item \textbf{Solver (cloud side)}: sweep $\rho$ and $\kappa(PS)$, and two-level and GMG iterations with the sweep
        and element-Schwarz smoothers, on:
        \begin{itemize}[nosep]
          \item production L0 and L1 with the fields on;
          \item the zero-field deep cathode mesh;
          \item $p=2$ and $p=4$.
        \end{itemize}
        Compare variant A, variant B and the current norm.
  \item \textbf{Accuracy (either side)}: single-slab errors against the current norm and the full graph norm, in
        particular on the $x$-refined test where $C=1.6$ was found dissipative.
  \item \textbf{Coupled stability (local side)}: the cathode and anode mode tests.
\end{enumerate}

\paragraph{Where the measurements come from.}  {\sloppy The table is from the cloud branch
\texttt{separable-gram-fast-diagonalization}.  The code is in \texttt{MatrixFreeDPGOperator}: \texttt{transportSweep()},
and \texttt{upwindCellOrder()} with weighted cycle breaking.  The driver options are:
\begin{itemize}[nosep]
  \item \texttt{--matrixFreeMultiplicativeSchwarz}
  \item \texttt{--matrixFreeSweepSmoother}
  \item \texttt{--matrixFreeWeightedCycleBreaking}
  \item \texttt{--matrixFreeCheckAfterSolve}
\end{itemize}
All runs are at $p=2$, $\Delta p=2$, on one core.  The L1 runs with fields on are in progress.\par}

\end{document}
